\section{Minimal Theoretical Model}
\label{sec:minimal_model}

The purpose of this section is to construct the simplest dynamical model capable of representing nuclear-state control of fission. The model is intentionally abstract and does not initially assume a specific isotope.

Consider two nuclear configurations,
\begin{equation}
|A\rangle,
\qquad
|B\rangle,
\end{equation}
with different fission decay rates,
\begin{equation}
\Gamma_f^{(A)}
\neq
\Gamma_f^{(B)}.
\end{equation}

An external actuator couples the two states,
\begin{equation}
|A\rangle
\overset{u_Q(t)}{\longleftrightarrow}
|B\rangle.
\end{equation}

\subsection{Population-based two-state model}

Let $p_A(t)$ and $p_B(t)$ denote the two populations. The external actuator produces transition rates
\begin{equation}
W_{A\rightarrow B}(t),
\qquad
W_{B\rightarrow A}(t).
\end{equation}

Both states may decay through fission and through competing channels. Define
\begin{align}
\Gamma_A &= \Gamma_f^{(A)}+\Gamma_d^{(A)},\\
\Gamma_B &= \Gamma_f^{(B)}+\Gamma_d^{(B)}.
\end{align}

The population equations are
\begin{align}
\frac{dp_A}{dt}
&=
-
W_{A\rightarrow B}p_A
+
W_{B\rightarrow A}p_B
-
\Gamma_Ap_A,
\\
\frac{dp_B}{dt}
&=
W_{A\rightarrow B}p_A
-
W_{B\rightarrow A}p_B
-
\Gamma_Bp_B.
\end{align}

The instantaneous fission rate for the nuclear ensemble is
\begin{equation}
r_f(t)
=
\Gamma_f^{(A)}p_A(t)
+
\Gamma_f^{(B)}p_B(t),
\end{equation}
and the integrated fission probability is
\begin{equation}
P_f
=
\int_0^\infty r_f(t)\,dt.
\end{equation}

For an initially populated state $A$ and no external control,
\begin{equation}
P_f^{(0)}
=
\frac{\Gamma_f^{(A)}}{\Gamma_A}.
\end{equation}

If population can be transferred to state $B$, control is possible whenever the two branching ratios differ:
\begin{equation}
\boxed{
\frac{\Gamma_f^{(A)}}{\Gamma_A}
\neq
\frac{\Gamma_f^{(B)}}{\Gamma_B}
}
\end{equation}
together with
\begin{equation}
\boxed{
W_{A\rightarrow B}\neq0.
}
\end{equation}

No long-lived quantum coherence is required.

\subsection{Microscopic control efficiency}

We define a dimensionless microscopic control efficiency,
\begin{equation}
\eta_Q
=
\frac{P_f^{(Q)}-P_f^{(0)}}{P_f^{(0)}},
\end{equation}
where $P_f^{(Q)}$ denotes the fission probability under external actuation.

Positive $\eta_Q$ corresponds to enhanced fission, negative $\eta_Q$ to suppression. The quantity provides a direct metric that can be calculated microscopically or measured experimentally.

\subsection{Competition between control and relaxation}

State transfer must compete with intrinsic relaxation and equilibration. Introduce
\begin{equation}
\tau_Q\sim\frac{1}{W_Q}
\end{equation}
and
\begin{equation}
\tau_{\mathrm{eq}}
\sim
\frac{1}{\Gamma_{\mathrm{eq}}}.
\end{equation}

A necessary condition for the prepared state to influence fission is approximately
\begin{equation}
\tau_Q
\lesssim
\tau_{\mathrm{eq}},
\end{equation}
or
\begin{equation}
W_Q
\gtrsim
\Gamma_{\mathrm{eq}}.
\end{equation}

We therefore define
\begin{equation}
\mathcal{C}_Q
=
\frac{W_Q}{\Gamma_{\mathrm{eq}}},
\end{equation}
with $\mathcal{C}_Q\ll1$, $\sim1$, and $\gg1$ representing ineffective, competing, and dominant external control, respectively.

\subsection{Coherent two-state extension}

A coherent version of the model may be represented by
\begin{equation}
H(t)
=
\frac{\hbar}{2}
\begin{pmatrix}
-\Delta &
\Omega(t)e^{-i\phi(t)}
\\
\Omega(t)e^{i\phi(t)} &
\Delta
\end{pmatrix},
\end{equation}
where $\Omega(t)$ is a Rabi coupling, $\Delta$ the detuning, and $\phi(t)$ an externally controlled phase.

The density matrix evolves according to a master equation,
\begin{equation}
\frac{d\rho}{dt}
=
-\frac{i}{\hbar}[H,\rho]
+
\mathcal{L}_{\mathrm{decay}}[\rho]
+
\mathcal{L}_{\mathrm{decoh}}[\rho].
\end{equation}

The fission rate can then contain both population and interference contributions,
\begin{equation}
r_f(t)
=
\Gamma_f^{(A)}\rho_{AA}
+
\Gamma_f^{(B)}\rho_{BB}
+
r_f^{(\mathrm{int})}.
\end{equation}

If both states feed common fission channels,
\begin{equation}
r_f^{(\mathrm{int})}
\propto
2\,\mathrm{Re}
\left[
\rho_{AB}
A_{A\rightarrow f}
A_{B\rightarrow f}^{\ast}
\right].
\end{equation}

A response satisfying
\begin{equation}
\frac{\partial P_f}{\partial\phi}\neq0
\end{equation}
would constitute a particularly strong signature of genuine coherent fission control.

\subsection{Three regimes}

The model identifies three qualitative regimes.

\paragraph{Regime I: no nuclear-state control.}
If
\begin{equation}
W_Q\ll\Gamma_{\mathrm{eq}},
\end{equation}
the externally driven redistribution is negligible and
\begin{equation}
P_f^{(Q)}\simeq P_f^{(0)}.
\end{equation}

\paragraph{Regime II: population control.}
If externally driven redistribution competes with relaxation but coherence is rapidly lost, the density matrix is effectively diagonal and control proceeds through state populations.

\paragraph{Regime III: coherent control.}
If the coherent coupling satisfies approximately
\begin{equation}
\Omega\gtrsim\Gamma_{\mathrm{decoh}}
\end{equation}
over the relevant interval, phase-sensitive interference may become observable.

The second regime is sufficient for the general reactor-control framework and is considerably less demanding than the third.

\subsection{Coupling the microscopic model to a neutron field}

For a two-state population model,
\begin{equation}
\sigma_f^{\mathrm{eff}}(E,t)
=
p_A(t)\sigma_f^{(A)}(E)
+
p_B(t)\sigma_f^{(B)}(E).
\end{equation}

The macroscopic fission rate becomes
\begin{equation}
R_f(t)
=
N_f
\int
\Phi(E,t)
\left[
p_A(t)\sigma_f^{(A)}(E)
+
p_B(t)\sigma_f^{(B)}(E)
\right]
dE.
\label{eq:macro_two_state}
\end{equation}

Thus,
\begin{equation}
u_Q(t)
\rightarrow
p_i(t)
\rightarrow
\sigma_f^{\mathrm{eff}}(t)
\rightarrow
R_f(t).
\end{equation}

\subsection{Effective reactivity contribution}

Suppose external control induces
\begin{equation}
\sigma_f
\rightarrow
\sigma_f+\delta\sigma_f.
\end{equation}

Then
\begin{equation}
\delta\Sigma_f
=
N_f\delta\sigma_f.
\end{equation}

For a sufficiently small perturbation,
\begin{equation}
\delta k_{\mathrm{eff}}
\simeq
\frac{\partial k_{\mathrm{eff}}}{\partial\Sigma_f}
\delta\Sigma_f.
\end{equation}

Since
\begin{equation}
\rho
=
1-\frac{1}{k_{\mathrm{eff}}},
\end{equation}
the associated first-order reactivity change is
\begin{equation}
\delta\rho_Q
\simeq
\frac{\delta k_{\mathrm{eff}}}{k_{\mathrm{eff}}^2}.
\end{equation}

Near criticality,
\begin{equation}
k_{\mathrm{eff}}\simeq1,
\end{equation}
and therefore
\begin{equation}
\delta\rho_Q\simeq\delta k_{\mathrm{eff}}.
\end{equation}

\subsection{Extension of point-reactor kinetics}

Ordinary point kinetics may be written as
\begin{equation}
\frac{dn}{dt}
=
\frac{\rho-\beta}{\Lambda}n
+
\sum_i\lambda_i C_i.
\end{equation}

In the present framework,
\begin{equation}
\rho
=
\rho_{\mathrm{conv}}
+
\rho_Q,
\end{equation}
so that
\begin{equation}
\frac{dn}{dt}
=
\frac{
\rho_{\mathrm{conv}}
+
\rho_Q
-
\beta
}{
\Lambda
}n
+
\sum_i\lambda_iC_i.
\label{eq:extended_point_kinetics}
\end{equation}

Unlike an externally prescribed conventional reactivity insertion, $\rho_Q$ has its own microscopic dynamics:
\begin{align}
\dot{\rho}_N
&=
\mathcal{F}(\rho_N,u_Q,\Phi,\ldots),
\\
\rho_Q
&=
\mathcal{G}(\rho_N),
\\
\dot n
&=
\frac{\rho_{\mathrm{conv}}+\rho_Q-\beta}{\Lambda}n
+
\sum_i\lambda_iC_i.
\end{align}

This system is the minimal mathematical architecture of a multiplying fission system whose fuel possesses an externally controllable internal nuclear-state variable.

\subsection{Feedback}

The coupling may be bidirectional. External actuation modifies the nuclear state,
\begin{equation}
u_Q\rightarrow\rho_N,
\end{equation}
while the neutron field itself can alter nuclear-state populations through absorption and compound-nucleus formation,
\begin{equation}
\Phi\rightarrow\rho_N.
\end{equation}
At the same time,
\begin{equation}
\rho_N
\rightarrow
\sigma_f^{\mathrm{eff}}
\rightarrow
\Phi.
\end{equation}

The resulting structure is therefore a feedback loop rather than a purely feed-forward actuator.

\subsection{Minimal falsifiable prediction}

Consider two otherwise identical experiments on a fissile target exposed to the same neutron distribution,
\begin{equation}
\Phi_1(E)=\Phi_2(E).
\end{equation}

If an external protocol prepares different nuclear-state distributions,
\begin{equation}
\rho_N^{(1)}
\neq
\rho_N^{(2)},
\end{equation}
the hypothesis predicts, under suitable conditions,
\begin{equation}
P_f^{(1)}
\neq
P_f^{(2)}.
\end{equation}

Equivalently,
\begin{equation}
\boxed{
\left.
\frac{\partial P_f}{\partial u_Q}
\right|_{\Phi}
\neq0.
}
\label{eq:falsifiable_prediction}
\end{equation}

Equation~\eqref{eq:falsifiable_prediction} is the simplest operational statement of the proposal. It separates state preparation from direct radiation-induced fission: the incident neutron field is held fixed while an independent actuator modifies the nuclear-state distribution.
